% Auto-generated by python/build_tables.py
\begin{table*}[!htbp]
\centering
\begin{threeparttable}
\caption{Readiness-threshold robustness.}
\label{tab:appendix-threshold-robustness}
\scriptsize
\begin{tabular}{rrrrrrr}
\toprule
\textbf{$\theta$} & \textbf{$q$} & \textbf{$BS_{low}$ conv.} & \textbf{$BS_{high}$ conv.} & \textbf{$BS_{low}$ $P_{95}$} & \textbf{$BS_{high}$ $P_{95}$} & \textbf{Reduction} \\
\midrule
0.75 & 0.80 & 100.0\% & 100.0\% & 6,582.6 & 4,221.5 & 35.9\% \\
0.80 & 0.70 & 100.0\% & 100.0\% & 9,762.5 & 5,070.0 & 48.1\% \\
0.80 & 0.80 & 100.0\% & 100.0\% & 12,794.1 & 5,781.3 & 54.8\% \\
0.80 & 0.90 & 100.0\% & 100.0\% & 19,046.1 & 6,921.5 & 63.7\% \\
0.85 & 0.80 & 27.0\% & 100.0\% & 50,000.0 & 10,436.4 & 79.1\% \\
\bottomrule
\end{tabular}
\begin{tablenotes}
\footnotesize
\item Notes: $\theta$ is the tie-level readiness threshold and $q$ is the boundary-level readiness threshold. The baseline condition is $\theta=0.80$ and $q=0.80$. The last row includes censoring at $T_{\max}=50{,}000$ because $BS_{low}$ frequently fails to reach readiness under the hardest tie-level threshold. Across all scenarios, translation-capable boundary spanning remains faster and more reliable than concentration without translation capability.
\end{tablenotes}
\end{threeparttable}
\end{table*}
